\documentclass[10pt,a4paper]{amsart}
\usepackage[margin=19mm]{geometry}
\usepackage{amsmath,amssymb,mathtools}
\usepackage[scr=rsfs,cal=euler]{mathalfa}
\usepackage{xfrac}
\usepackage[unicode,hidelinks,bookmarks=false]{hyperref}
\newcommand{\ie}{i.e.\ }
\newcommand{\cf}{cf.\ }
\newcommand{\resp}{respectively\ }
\newcommand{\Subsection}{\subsection}
\providecommand{\nobreakdash}{\nobreak\hbox{-}}
\setlength{\emergencystretch}{2em}
\multlinegap0pt
\usepackage{amssymb}
\usepackage{multirow}
\usepackage{xcolor}
\usepackage{booktabs}
\usepackage{algorithm}\usepackage{algpseudocode}
\usepackage{natbib}
\usepackage{bm}
\usepackage{esint}
\newtheorem{theorem}{Theorem}
\title{\textbf{Two-Dimensional $\beta$-plane Turbulence: Dual Cascade and Zonal Jets}}
\author{Yuri Cacchiò, Amirali Hannani, Gigliola Staffilani}
\date{Source: arXiv:2608.22021v1 (CC BY 4.0)}
\begin{document}
\maketitle

\noindent\emph{Source and licence.} \textbf{Two-Dimensional $\beta$-plane Turbulence: Dual Cascade and Zonal Jets}. Version arXiv:2608.22021v1 (CC BY 4.0), distributed by arXiv under the Creative Commons Attribution 4.0 International licence, http://creativecommons.org/licenses/by/4.0/, as recorded in the arXiv licence metadata for this exact version and shown on the abstract page https://arxiv.org/abs/2608.22021v1. Full licence notice: \texttt{licenses/arxiv-2608.22021-CC-BY-4.0.txt}.\par\smallskip

\noindent\textit{Editorial scope.} Counted unit: the article's theorem thm:vanishing\_coriolis (source0.tex lines 482--487). The article's complete original proof of it is delivered with it (source0.tex lines 622--704, one \texttt{proof} environment of 83 lines, which the article places away from the statement). No mathematics is written, repaired or re-derived: the statement and the proof are the article's own lines, in the article's own notation, and no retained line is substituted or re-flowed.  Retained uncounted as the source-local context the proof needs: lines 109--111; lines 377--379; lines 418--423; lines 442--446.  Nothing was omitted from any retained span, so the package carries no omission record.  No retained line contains a citation, so the wrapper carries no bibliography and no bibliography entry.  No retained line contains a figure environment, an image-inclusion command or drawing code, so no image is delivered and the package declares no asset dependency.  Editorial formatting only, in addition: an AMS \texttt{amsart} preamble that restates the article's own declarations and macros used by the retained lines, namely the environments \texttt{theorem} on separate counters, together with the standard packages those lines need.  The retained environments print in this document as Theorem 1 in order of appearance, whereas the article prints the same items with the numbering of its own full document.  The complete native source of the article as distributed in the arXiv e-print for this exact version is bundled unchanged as \texttt{sources/208295/source0.tex}.
\begin{equation}\label{eq:boundary_conditions}
    u(t, x^1+L, x^2) = u(t, x^1, x^2), \qquad u(t, x^1, a) = u(t, x^1, b) = 0.
\end{equation}

\begin{equation}\label{eq:abstract_NS}
    du + \big[ \nu A u + B(u,u) + \alpha u + \mathcal{C}(u) \big] dt = d\zeta(t),
\end{equation}

\begin{align}
    \Gamma(y) &= \mathbb{E} \fint_{\mathbb{T} \times \mathbb{R}} u(x) \otimes u(x+y) \, dx, \label{def:Gamma}\\
    D^j(y) &= \mathbb{E} \fint_{\mathbb{T} \times \mathbb{R}} (\delta_y u(x) \otimes \delta_y u(x)) \delta_y u^j(x) \, dx, \label{def:Dj}\\
    \varTheta(y) &= \frac{1}{2}\mathbb{E} \fint_{\mathbb{T} \times \mathbb{R}} \big( u(x) \otimes (f u^\perp)(x+y) + (f u^\perp)(x) \otimes u(x+y) \big) \, dx, \label{def:Theta}\\
    a(y) &= \frac{1}{2}\sum_{j} b_j^2 \fint_{\mathbb{T}\times \mathbb{R}} e_j(x) \otimes e_j(x+y) \, dx, \label{def:a}
\end{align}

\begin{align}
    \mathfrak{C}(y) &= \mathbb{E}\fint_{\mathbb{T}\times \mathbb{R}} \omega(x)\omega(x+y) \, dx, \label{def:vorticity_corr}\\
    \mathfrak{Q}(y) &= \frac{1}{2}\beta \, \mathbb{E}\fint_{\mathbb{T}\times \mathbb{R}} \big( u^2(x)\omega(x+y) + u^2(x+y)\omega(x) \big) \, dx, \label{def:coriolis_vort_corr}\\
    \mathfrak{a}(y) &= \frac{1}{2}\sum_{j=1}^\infty b_j^2\fint_{\mathbb{T}\times \mathbb{R}} \big(\nabla^\perp \cdot e_j(x)\big)\big(\nabla^\perp \cdot e_j(x+y)\big) \, dx. \label{def:noise_curl_corr}
\end{align}

\begin{theorem}[Vanishing of Coriolis Terms]\label{thm:vanishing_coriolis}
Let $u$ be a statistically stationary solution to the $\beta$-plane stochastic Navier-Stokes equations \eqref{eq:abstract_NS}. Then, the spherically averaged Coriolis contributions vanish for all scales $l > 0$:
\begin{equation}
    \overline{\varTheta}(l) = 0 \quad \text{and} \quad \overline{\mathfrak{Q}}(l) = 0.
\end{equation}
\end{theorem}

\begin{proof}[Proof of Theorem \ref{thm:vanishing_coriolis}]
We split the proof into two steps. 
Before proceeding, we note that the stream function $\psi$ is periodic in the zonal direction $x^1 \in \mathbb{T}$. Indeed, integrating the divergence-free condition over $\mathbb{T}$ yields
\begin{equation}
    \partial_2 \int_{\mathbb{T}} u^2 dx^1 = 0,
\end{equation}
that is,
\begin{equation}\label{eq:integral}
    \int_{\mathbb{T}} u^2 dx^1 = C
\end{equation}
for $x^2\in[a,b]$. Since $u^2 = 0$ at the boundary $x^2=a$, the integral \eqref{eq:integral} vanishes everywhere, ensuring
\begin{equation}
    \psi(L, x^2) - \psi(0, x^2) = \int_{\mathbb{T}} u^2 dx^1 = 0.
\end{equation}

\noindent\textbf{Step 1: Vanishing of the Coriolis enstrophy flux $\mathfrak{Q}(y)=0$.}\\
Restricting the domain of integration to the periodic channel, the enstrophy Coriolis correlation \eqref{def:coriolis_vort_corr} reads 
\begin{equation}\label{eq:Q_restricted}
    \mathfrak{Q}(y) = \frac{1}{2}\beta \, \mathbb{E} \fint_{\mathbb{T} \times [a, b]} \big( u^2(x) \omega(x+y) + \omega(x) u^2(x+y) \big) \, dx.
\end{equation}
Substituting $\omega = \partial_{11}^2 \psi + \partial_{22}^2 \psi$, the integral decomposes into a zonal component $Z(y)$ and a meridional component $M(y)$:
\begin{align*}
    \fint_{\mathbb{T} \times [a, b]} \big[ \partial_1 \psi(x) \Delta \psi(x+y) + \Delta \psi&(x) \partial_1 \psi(x+y) \big] \, dx\\
    =&\fint_{\mathbb{T} \times [a, b]} \big[ \partial_1 \psi(x) \partial_{11}^2 \psi(x+y) + \partial_{11}^2 \psi(x) \partial_1 \psi(x+y) \big] \, dx\\
    +&\fint_{\mathbb{T} \times [a, b]} \big[ \partial_1 \psi(x) \partial_{22}^2 \psi(x+y) + \partial_{22}^2 \psi(x) \partial_1 \psi(x+y) \big] \, dx\\
    =: &Z(y) + M(y).
\end{align*}
For $Z(y)$, integrating the second term by parts with respect to $x^1$ shifts one derivative from $x$ to $x+y$. Because $\psi$ is periodic in $x^1$, the boundary terms vanish, yielding $Z(y) = 0$.

For $M(y)$, we integrate both terms by parts with respect to $x^2 \in [a, b]$,
\begin{align*}
    \fint_a^b \partial_1 \psi(x) \partial_{22}^2 \psi(x+y) \, dx^2 &= \Big[ \partial_1 \psi(x) \partial_2 \psi(x+y)\Big]_{a}^{b} - \fint_a^b \partial_{12}^2 \psi(x) \partial_2 \psi(x+y) \, dx^2, \\
    \fint_a^b \partial_{22}^2 \psi(x) \partial_1 \psi(x+y) \, dx^2 &= \Big[ \partial_2 \psi(x) \partial_1 \psi(x+y) \Big]_{a}^{b} - \fint_a^b \partial_{2} \psi(x) \partial_{21}^2 \psi(x+y) \, dx^2.
\end{align*}
The no-slip boundary condition \eqref{eq:boundary_conditions} ensures $u^2 = \partial_1 \psi = 0$ and $u^1 = -\partial_2 \psi = 0$ for $x^2 \in \{a, b\}$, canceling the two brackets. We are left with
\begin{equation*}
    M(y) = \fint_{\mathbb{T} \times I} \big[ - \partial_{12}^2 \psi(x) \partial_2 \psi(x+y) - \partial_2 \psi(x) \partial_{21}^2 \psi(x+y) \big] \, dx.
\end{equation*}
Integrating the first term by parts along the periodic zonal direction $x^1$, we get 
\begin{equation*}
    - \fint_a^b dx^2 \fint_{\mathbb{T}} \partial_1 \big(\partial_2 \psi(x)\big) \partial_2 \psi(x+y) \, dx^1 = \fint_{\mathbb{T} \times [a, b]} \partial_2 \psi(x) \partial_{12}^2 \psi(x+y) \, dx,
\end{equation*}
where the boundary term vanishes by periodicity. Since the $H^3$ regularity of the stream function $\psi$ guarantees the commutation of the derivatives in the weak sense $\partial_{12}^2 \psi = \partial_{21}^2 \psi$ in $L^2$, this remaining term cancels the second term of $M(y)$, yielding $M(y) = 0$. Therefore $\mathfrak{Q}(y)=0$ for all $y\in \mathbb{R}^2$, and spherical averaging yields $\overline{\mathfrak{Q}}(l)=0$. 

\text{}




\noindent \textbf{Step 2: Vanishing of the velocity Coriolis trace $\text{Tr}(\varTheta(y))=0$.}\\
Recalling $u^\perp = (-u^2, u^1)$ and the definition \eqref{def:Theta}, we write
\begin{equation}
    \text{Tr}(\varTheta(y)) = \frac{1}{2}\mathbb{E} \fint_{\mathbb{T}\times [a,b]} \big(u(x) \cdot (fu^\perp)(x+y) + (f u^\perp)(x) \cdot u(x+y)\big)\, dx.
\end{equation}
Evaluating the shift $f(x+y) - f(x) = \beta y^2$, the uniform rotation component $f_0$ cancels, and we obtain
\begin{align}
    \text{Tr}(\varTheta(y)) &= \frac{1}{2}\beta y^2 \, \mathbb{E} \fint_{\mathbb{T} \times I} \big( u^2(x)u^1(x+y) - u^1(x)u^2(x+y) \big) \, dx\notag \\
    &= \frac{1}{2}\beta y^2 \, \mathbb{E} \fint_{\mathbb{T} \times I} \big( -\partial_1 \psi(x) \partial_2 \psi(x+y) + \partial_2 \psi(x) \partial_1 \psi(x+y) \big) \, dx.\label{eq:trace_psi}
\end{align}
The no-slip boundary condition \eqref{eq:boundary_conditions} implies $\psi$ is constant for $x^2\in \{a,b\}$. We set $\psi(x^1, a) = 0$. On $x^2=b$, the constant is given by the net zonal transport,
\begin{equation}
    \psi(x^1, b) = - \fint_a^b u^1 dx^2 := C_b.
\end{equation}
Integrating the first term of \eqref{eq:trace_psi} by parts with respect to $x^1$, we derive
\begin{equation}\label{eq:int1}
    \fint_a^b dx^2 \fint_{\mathbb{T}} -\partial_1 \psi (x)\partial_2 \psi(x+y) \, dx^1 = \fint_{\mathbb{T} \times [a, b]} \psi(x) \, \partial_{12}^2 \psi(x+y) \, dx,
\end{equation}
where again by periodicity the boundary term vanishes. Integrating the second term of \eqref{eq:trace_psi} by parts with respect to $x^2$, we get
\begin{align}
    \fint_{\mathbb{T}} dx^1 \fint_a^b \partial_2 \psi(x) \partial_1\psi&(x+y) \, dx^2 \label{eq:int2}\\
    &= \fint_{\mathbb{T}} dx^1 \left( - \fint_a^b \psi(x) \, \partial_{21}^2 \psi(x+y) \, dx^2 + \Big[ \psi(x) \, \partial_1 \psi(x+y) \Big]_{x^2=a}^{x^2=b} \right)\notag \\
    &= - \fint_{\mathbb{T} \times [a, b]} \psi(x) \, \partial_{21}^2 \psi(x+y) \, dx + C_b \fint_{\mathbb{T}} \partial_1 \psi(x^1+y^1, b+y^2) \, dx^1.\notag
\end{align}
Notice that the boundary term evaluated at $x^2=b$ is the integral of a derivative over the periodic domain $\mathbb{T}$. Hence,
\begin{equation*}
    C_b \fint_{\mathbb{T}} \partial_1 \psi(x^1+y^1, b+y^2) \, dx^1 = C_b \big[ \psi(L+y^1, b+y^2) - \psi(y^1, b+y^2) \big] = 0.
\end{equation*}
Summing the two integrated components \eqref{eq:int1} and \eqref{eq:int2}, all boundary terms vanish independently of the constant $C_b$, yielding 
\begin{equation}\label{eq:coclusionvanish}
    \fint_{\mathbb{T} \times [a, b]} \psi(x) \big[ \partial_{12}^2 - \partial_{21}^2 \big] \psi(x+y) \, dx = 0.
\end{equation}
Therefore, $\text{Tr}(\varTheta(y)) = 0$ for all $y \in \mathbb{R}^2$, implying $\overline{\varTheta}(l) = 0$.
\end{proof}

\end{document}
